% GENERATED FILE. Edit canonical lessons, then run npm run generate:books.
% canonical-source-hash: fnv1a64:cfb515f887b02c98
% canonical-lessons: RU-C479-doveriye, RU-C479-doveryat, RU-C479-predavat, RU-C479-uvazhat, RU-C479-iskrenniy

\chapter{Trust and Betrayal}
\label{ch:ru-trust-and-betrayal}
\hlchaptermodality{479}

\begin{chapteropening}
\textbf{By the end of this chapter:} \emph{I can talk about trust, respect and betrayal.}

Complete the last lesson of chapter 479: I can talk about trust, respect and betrayal.
\end{chapteropening}

\section[\texorpdfstring{\ru{доверие} (dovériye)}{доверие (dovériye)}]{\ru{доверие} (dovériye) — trust}

\label{lesson:RU-C479-doveriye}



\begin{quote}
Before the new one: say the Russian for to play up, then the Russian for to grumble.
\end{quote}

\subsection*{You'll want to know: \ru{доверие}}
\textbf{\ru{доверие}} (\emph{dovériye}) — \textquotedblleft{}trust\textquotedblright{}.

\begin{quote}
\ru{Между} \ru{нами} \ru{полное} \ru{доверие}.
\end{quote}

\begin{quote}
\emph{Mézhdu námi pólnoye dovériye.}
\end{quote}

\begin{quote}
There is complete trust between us.
\end{quote}

Neuter.

\begin{cousinweb}
\emph{\ru{доверие}} is built on \emph{\ru{до}-} and \emph{\ru{вера}} 'faith', which is inherited from Proto-Slavic \emph{věra} 'faith, belief' (EtymDB).
\end{cousinweb}

\begin{grammarlens}[title={What you've built}]
One more word for telling, explaining and describing at B1.
\end{grammarlens}

\subsection*{Your turn}
\begin{itemize}
\raggedright
  \item \emph{Say it:} \emph{dovériye}
  \item \emph{Say it:} \emph{dovériye}, once more
  \item \emph{Say it:} a sentence of your own with \emph{dovériye}
  \item \emph{Recall:} say the Russian for to play up, then the Russian for to grumble, then say \emph{dovériye} again
\end{itemize}

\begin{tcolorbox}[breakable,colback=teal!4,colframe=teal!35!black,title={Before you move on}]
What does \ru{доверие} mean? (\textquotedblleft{}Trust\textquotedblright{}.) Say it once more.
\end{tcolorbox}
\section[\texorpdfstring{\ru{доверять} (doveryát)}{доверять (doveryát)}]{\ru{доверять} (doveryát) — to trust}

\label{lesson:RU-C479-doveryat}



\begin{quote}
Before the new one: say the Russian for to grumble, then the Russian for trust.
\end{quote}

\subsection*{You'll want to know: \ru{доверять}}
\textbf{\ru{доверять}} (\emph{doveryát}) — \textquotedblleft{}to trust\textquotedblright{}.

\begin{quote}
\ru{Я} \ru{доверяю} \ru{своему} \ru{врачу}.
\end{quote}

\begin{quote}
\emph{Ya doveryáyu svoyemú vrachú.}
\end{quote}

\begin{quote}
I trust my doctor.
\end{quote}

Imperfective; perfective \ru{доверить}. Takes the dative: \emph{\ru{доверять} \ru{врачу}}.

\begin{cousinweb}
\emph{\ru{доверять}} is built on \emph{\ru{до}-} and the root of \emph{\ru{вера}} 'faith', inherited from Proto-Slavic \emph{věra} (EtymDB).
\end{cousinweb}

\begin{grammarlens}[title={What you've built}]
One more word for telling, explaining and describing at B1.
\end{grammarlens}

\subsection*{Your turn}
\begin{itemize}
\raggedright
  \item \emph{Say it:} \emph{doveryát}
  \item \emph{Say it:} \emph{doveryát}, once more
  \item \emph{Say it:} a sentence of your own with \emph{doveryát}
  \item \emph{Recall:} say the Russian for to grumble, then the Russian for trust, then say \emph{doveryát} again
\end{itemize}

\begin{tcolorbox}[breakable,colback=teal!4,colframe=teal!35!black,title={Before you move on}]
What does \ru{доверять} mean? (\textquotedblleft{}To trust\textquotedblright{}.) Say it once more.
\end{tcolorbox}
\section[\texorpdfstring{\ru{предавать} (predavát)}{предавать (predavát)}]{\ru{предавать} (predavát) — to betray}

\label{lesson:RU-C479-predavat}



\begin{quote}
Before the new one: say the Russian for trust, then the Russian for to trust.
\end{quote}

\subsection*{You'll want to know: \ru{предавать}}
\textbf{\ru{предавать}} (\emph{predavát}) — \textquotedblleft{}to betray\textquotedblright{}.

\begin{quote}
\ru{Настоящий} \ru{друг} \ru{никогда} \ru{не} \ru{предаёт}.
\end{quote}

\begin{quote}
\emph{Nastoyáshchiy drug nikogdá ne predayót.}
\end{quote}

\begin{quote}
A true friend never betrays you.
\end{quote}

Imperfective; present \ru{предаю}, \ru{предаёшь} (-\ru{ава}- becomes -\ru{а}-). Perfective \ru{предать}.

\begin{cousinweb}
\emph{\ru{предавать}} is built on \emph{\ru{пред}-} and \emph{\ru{давать}} 'to give', so to betray is to hand someone over; the root is inherited from Proto-Slavic \emph{dati} 'to give' (EtymDB).
\end{cousinweb}

\begin{grammarlens}[title={What you've built}]
One more word for telling, explaining and describing at B1.
\end{grammarlens}

\subsection*{Your turn}
\begin{itemize}
\raggedright
  \item \emph{Say it:} \emph{predavát}
  \item \emph{Say it:} \emph{predavát}, once more
  \item \emph{Say it:} a sentence of your own with \emph{predavát}
  \item \emph{Recall:} say the Russian for trust, then the Russian for to trust, then say \emph{predavát} again
\end{itemize}

\begin{tcolorbox}[breakable,colback=teal!4,colframe=teal!35!black,title={Before you move on}]
What does \ru{предавать} mean? (\textquotedblleft{}To betray\textquotedblright{}.) Say it once more.
\end{tcolorbox}
\section[\texorpdfstring{\ru{уважать} (uvazhát)}{уважать (uvazhát)}]{\ru{уважать} (uvazhát) — to respect}

\label{lesson:RU-C479-uvazhat}



\begin{quote}
Before the new one: say the Russian for to trust, then the Russian for to betray.
\end{quote}

\subsection*{You'll want to know: \ru{уважать}}
\textbf{\ru{уважать}} (\emph{uvazhát}) — \textquotedblleft{}to respect\textquotedblright{}.

\begin{quote}
\ru{Студенты} \ru{уважают} \ru{этого} \ru{профессора}.
\end{quote}

\begin{quote}
\emph{Studénty uvazháyut étogo proféssora.}
\end{quote}

\begin{quote}
The students respect this professor.
\end{quote}

Imperfective. The noun is \emph{\ru{уважение}} — respect.

\begin{cousinweb}
\emph{\ru{уважать}} was borrowed from Polish \emph{uważać} 'to consider' (EtymDB).
\end{cousinweb}

\begin{grammarlens}[title={What you've built}]
One more word for telling, explaining and describing at B1.
\end{grammarlens}

\subsection*{Your turn}
\begin{itemize}
\raggedright
  \item \emph{Say it:} \emph{uvazhát}
  \item \emph{Say it:} \emph{uvazhát}, once more
  \item \emph{Say it:} a sentence of your own with \emph{uvazhát}
  \item \emph{Recall:} say the Russian for to trust, then the Russian for to betray, then say \emph{uvazhát} again
\end{itemize}

\begin{tcolorbox}[breakable,colback=teal!4,colframe=teal!35!black,title={Before you move on}]
What does \ru{уважать} mean? (\textquotedblleft{}To respect\textquotedblright{}.) Say it once more.
\end{tcolorbox}
\section[\texorpdfstring{\ru{искренний} (ískrenniy)}{искренний (ískrenniy)}]{\ru{искренний} (ískrenniy) — sincere}

\label{lesson:RU-C479-iskrenniy}



\begin{quote}
Before the new one: say the Russian for to betray, then the Russian for to respect.
\end{quote}

\subsection*{You'll want to know: \ru{искренний}}
\textbf{\ru{искренний}} (\emph{ískrenniy}) — \textquotedblleft{}sincere\textquotedblright{}.

\begin{quote}
\ru{Она} \ru{дала} \ru{мне} \ru{искренний} \ru{совет}.
\end{quote}

\begin{quote}
\emph{Oná dalá mne ískrenniy sovét.}
\end{quote}

\begin{quote}
She gave me sincere advice.
\end{quote}

Soft adjective: \ru{искренняя}, \ru{искреннее}. The adverb is \emph{\ru{искренне}} — sincerely.

\begin{cousinweb}
\emph{\ru{искренний}} is inherited from Proto-Slavic \emph{iskr\ru{ь}nj\ru{ь}} (EtymDB), usually glossed 'near, close'.
\end{cousinweb}

\begin{grammarlens}[title={What you've built}]
One more word for telling, explaining and describing at B1.
\end{grammarlens}

\subsection*{Your turn}
\begin{itemize}
\raggedright
  \item \emph{Say it:} \emph{ískrenniy}
  \item \emph{Say it:} \emph{ískrenniy}, once more
  \item \emph{Say it:} a sentence of your own with \emph{ískrenniy}
  \item \emph{Recall:} say the Russian for to betray, then the Russian for to respect, then say \emph{ískrenniy} again
\end{itemize}

\begin{tcolorbox}[breakable,colback=teal!4,colframe=teal!35!black,title={Before you move on}]
What does \ru{искренний} mean? (\textquotedblleft{}Sincere\textquotedblright{}.) Say it once more.
\end{tcolorbox}
